The affine cocycles of the preceding section have an infinitesimal counterpart.
Translation and dilation generate the Lie algebra of the positive affine group, and
their linear combinations give the continuous propagation law.  The metric
formulation introduced below is intrinsic: the symbols \(u\) and \(v\) label
arithmetic directions, but need not be coordinates.

\subsection{The affine Lie algebra and its flow}

On the assignment line with coordinate \(z\), let
\[
 E=\partial_z,
 \qquad
 H=z\partial_z .
\]
These are respectively the infinitesimal generators of translation and positive
dilation, and they satisfy \([E,H]=E\).  Fix constants
\(\mu,\lambda\in\mathbb R\), with \(\mu\ne0\).  The generator selected by an angle
\(\theta\) is
\begin{equation}\label{eq:angular-affine-generator}
 \Omega_\theta
 =\mu\cos\theta\,E+\lambda\sin\theta\,H .
\end{equation}
Consequently, an integral curve \(z=z(s)\) satisfies
\begin{equation}\label{eq:continuous-affine-flow}
 \frac{dz}{ds}
 =\mu\cos\theta+\lambda z\sin\theta .
\end{equation}
This Lie-algebra calculation is the primary derivation of the continuous AEG
equation.  It also fixes the convention: the translation coefficient is
\(\mu\cos\theta\), while the dilation coefficient is
\(\lambda\sin\theta\).

When \(\theta\) is constant, set \(b=\lambda\sin\theta\).  The solution through
\(z(0)=z_0\) is
\begin{equation}\label{eq:directional-affine-solution}
 z(s)=
 \begin{cases}
  e^{bs}z_0+\displaystyle\mu\cos\theta\,\frac{e^{bs}-1}{b},
     & b\ne0,\\[7pt]
  z_0+\mu s\cos\theta, & b=0.
 \end{cases}
\end{equation}
In particular, the positive arithmetic axes give
\[
 \theta=0:\quad z(s)=z_0+\mu s,
 \qquad
 \theta=\frac{\pi}{2}:\quad z(s)=e^{\lambda s}z_0.
\]
The opposite axes give subtraction and inverse dilation.

\subsection{The Pfaffian propagation constraint}

There is a useful three-dimensional local state space with coordinates
\((u,v,z)\).  On it define
\begin{equation}\label{eq:affine-pfaffian-form}
 \alpha=dz-\mu\,du-\lambda z\,dv
\end{equation}
and the two horizontal vector fields
\begin{equation}\label{eq:affine-horizontal-fields}
 D_u=\partial_u+\mu\partial_z,
 \qquad
 D_v=\partial_v+\lambda z\partial_z.
\end{equation}
A curve tangent to
\(\cos\theta\,D_u+\sin\theta\,D_v\) satisfies
\eqref{eq:continuous-affine-flow}.  Thus the frequently used notation
\[
 dz=\mu\,du+\lambda z\,dv
\]
means that \(\alpha\) vanishes on admissible tangent vectors.  It is not an
identity asserting that \(z\) is a function of two independent coordinates whose
ordinary partial derivatives are simultaneously \(\mu\) and \(\lambda z\).
Indeed,
\begin{equation}\label{eq:horizontal-noncommutation}
 [D_u,D_v]=\mu\lambda\,\partial_z.
\end{equation}
When \(\mu\lambda\ne0\), this bracket already rules out such an integrable
coordinate interpretation.

\subsection{Regular arithmetic expression spaces}

\begin{definition}[Regular arithmetic expression space]\label{def:regular-aes}
A \emph{regular arithmetic expression space}, or regular AES, is a tuple
\[
 \mathfrak E=(M,g,a;\mu,\lambda)
\]
such that \(M\) is an oriented smooth surface, possibly with boundary, \(g\) is a smooth Riemannian
metric, \(a\in C^\infty(M,\mathbb R)\), \(\mu\in\mathbb R\setminus\{0\}\),
\(\lambda\in\mathbb R\), and
\begin{equation}\label{eq:regular-aes-eikonal}
 |\nabla a|_g^2=\mu^2+\lambda^2a^2
\end{equation}
everywhere on \(M\).  The function \(a\) is called the
\emph{assignment function}.
\end{definition}

The metric in Definition~\ref{def:regular-aes} is part of the structure.
Equation~\eqref{eq:regular-aes-eikonal} is therefore an eikonal equation on a
specified Riemannian surface, not a metric-free identity.

Let \(J\) denote rotation by \(+\pi/2\) in the oriented tangent planes.  Since
\(\mu\ne0\), the quantity
\[
 q=\sqrt{\mu^2+\lambda^2a^2}
\]
is everywhere positive.  The next proposition shows that the eikonal definition
contains exactly the oriented orthonormal directional data used in the flow
equation.

\begin{proposition}[Canonical arithmetic frame]\label{prop:canonical-arithmetic-frame}
For an oriented Riemannian surface \((M,g)\), a smooth function \(a\), and
constants \(\mu\ne0\) and \(\lambda\), the following are equivalent.
\begin{enumerate}
\item The eikonal identity \eqref{eq:regular-aes-eikonal} holds.
\item There is a unique positively oriented orthonormal frame
      \((X_u,X_v)\) satisfying
      \begin{equation}\label{eq:arithmetic-frame-derivatives}
       X_u a=\mu,
       \qquad
       X_v a=\lambda a.
      \end{equation}
\end{enumerate}
When these conditions hold, put \(E_a=\nabla a/q\).  The frame is
\begin{equation}\label{eq:explicit-arithmetic-frame}
 \boxed{
 \begin{aligned}
 X_u&=\frac{\mu}{q}E_a-\frac{\lambda a}{q}J E_a,\\
 X_v&=\frac{\lambda a}{q}E_a+\frac{\mu}{q}J E_a.
 \end{aligned}}
\end{equation}
\end{proposition}

\begin{proof}
Assume the eikonal identity.  Then \(|\nabla a|=q>0\), so \(E_a\) and
\(JE_a\) form a smooth positively oriented orthonormal frame.  The coefficient
matrix in \eqref{eq:explicit-arithmetic-frame} is
\[
 \frac1q
 \begin{pmatrix}
  \mu&\lambda a\\
  -\lambda a&\mu
 \end{pmatrix},
\]
whose columns are orthonormal and whose determinant is one.  Hence
\((X_u,X_v)\) is positively oriented and orthonormal.  Since
\(da(E_a)=q\) and \(da(JE_a)=0\), the two identities in
\eqref{eq:arithmetic-frame-derivatives} follow.

For uniqueness, write any positively oriented orthonormal frame in the form
\[
 X_u=cE_a+sJE_a,
 \qquad
 X_v=-sE_a+cJE_a,
 \qquad c^2+s^2=1.
\]
The two directional identities force \(qc=\mu\) and \(-qs=\lambda a\).
Thus \(c=\mu/q\) and \(s=-\lambda a/q\), which gives
\eqref{eq:explicit-arithmetic-frame}.

Conversely, if an orthonormal frame satisfies
\eqref{eq:arithmetic-frame-derivatives}, then
\[
 |\nabla a|^2=(X_u a)^2+(X_v a)^2
 =\mu^2+\lambda^2a^2.
\]
\end{proof}

The frame is generally non-coordinate.  Applying its bracket to the assignment
function gives
\begin{equation}\label{eq:arithmetic-frame-bracket-on-a}
 [X_u,X_v]a
 =X_u(\lambda a)-X_v(\mu)
 =\mu\lambda.
\end{equation}
Thus, when \(\lambda\ne0\), the fields cannot commute.  If
\((\eta^u,\eta^v)\) is their dual coframe, the valid intrinsic Pfaffian statement is
\[
 da=\mu\,\eta^u+\lambda a\,\eta^v;
\]
the one-forms \(\eta^u,\eta^v\) are not presumed exact.

\begin{theorem}[Continuous affine flow]\label{thm:continuous-affine-flow}
Let \(\mathfrak E=(M,g,a;\mu,\lambda)\) be a regular AES with its canonical
arithmetic frame.  If a unit-speed curve \(\gamma\) has tangent
\[
 \dot\gamma
 =\cos\theta\,X_u+\sin\theta\,X_v,
\]
where \(\theta\) may vary along the curve, then
\begin{equation}\label{eq:intrinsic-continuous-affine-flow}
 \frac{d}{ds}(a\circ\gamma)
 =\mu\cos\theta+\lambda(a\circ\gamma)\sin\theta.
\end{equation}
For constant \(\theta\), its solutions are given by
\eqref{eq:directional-affine-solution}.
\end{theorem}

\begin{proof}
By the chain rule and \eqref{eq:arithmetic-frame-derivatives},
\[
 \frac{d}{ds}(a\circ\gamma)
 =da(\dot\gamma)
 =\cos\theta\,X_u a+\sin\theta\,X_v a
 =\mu\cos\theta+\lambda a\sin\theta.
\]
For constant \(\theta\), this is the linear equation already integrated in
\eqref{eq:directional-affine-solution}.
\end{proof}

\subsection{Rectification and the local order defect}

For \(\lambda\ne0\), define
\begin{equation}\label{eq:rectifying-assignment}
 r=\operatorname{arcsinh}\!\left(\frac{\lambda a}{\mu}\right).
\end{equation}
The chain rule and \eqref{eq:regular-aes-eikonal} give
\begin{equation}\label{eq:rectified-eikonal}
 |\nabla r|=|\lambda|.
\end{equation}
Thus \(r\) has constant gradient magnitude.  When \(\lambda=0\), the original
equation already reduces to \(|\nabla a|=|\mu|\), and \(a/\mu\) is a
unit-gradient rectifying function.

The same affine convention fixes the sign of the elementary local order defect.
For small increments \(h\) and \(k\), let
\[
 A_h(z)=z+\mu h,
 \qquad
 M_k(z)=e^{\lambda k}z.
\]
Then
\begin{align}
 (M_k\circ A_h)(z)-(A_h\circ M_k)(z)
 &=\mu h\bigl(e^{\lambda k}-1\bigr)\notag\\
 &=\mu\lambda\,hk+O\!\left(\lVert(h,k)\rVert^3\right).
 \label{eq:infinitesimal-order-defect}
\end{align}
This is an infinitesimal two-order comparison, not an exact area formula for an
arbitrary finite region.

\subsection{Affine flow inside projective flow}

An infinitesimal projective transformation of the line has the Riccati form
\[
 \dot z=\beta+\alpha z+\kappa z^2.
\]
The theory developed in this section is the affine slice \(\kappa=0\), with
\(\beta=\mu\cos\theta\) and \(\alpha=\lambda\sin\theta\).  Results proved for
regular affine AESs below do not automatically extend to the quadratic projective
flow.
